\documentclass[conference]{IEEEtran}
\IEEEoverridecommandlockouts
% The preceding line is only needed to identify funding in the first footnote. If that is unneeded, please comment it out.
\usepackage{cite}
\usepackage{amsmath,amssymb,amsfonts}
\usepackage{algorithmic}
\usepackage{graphicx}
\usepackage{textcomp}
\usepackage{xcolor}
\usepackage{url}
\usepackage{booktabs}
\usepackage{balance}

\def\BibTeX{{\rm B\kern-.05em{\sc i\kern-.025em b}\kern-.08em
    T\kern-.1667em\lower.7ex\hbox{E}\kern-.125emX}}

\begin{document}

\title{ThreatWeave: An Agentic Multimodal Framework for Cross-Modal Digital Communication Threat Analysis}

\author{
\IEEEauthorblockN{Ravulapalli Dharma Teja}
\IEEEauthorblockA{\textit{Dept. of Computer Science and Engineering} \\
\textit{Godavari Institute of Engineering and Technology (Autonomous)}\\
Rajahmundry, Andhra Pradesh, India \\
23551A05D1@giet.ac.in}
\and
\IEEEauthorblockN{Kotipalli Hemanth}
\IEEEauthorblockA{\textit{Dept. of Computer Science and Engineering} \\
\textit{Godavari Institute of Engineering and Technology (Autonomous)}\\
Rajahmundry, Andhra Pradesh, India \\
24555A0508@giet.ac.in}
\and
\IEEEauthorblockN{Gaddigopula Prabash}
\IEEEauthorblockA{\textit{Dept. of Computer Science and Engineering} \\
\textit{Godavari Institute of Engineering and Technology (Autonomous)}\\
Rajahmundry, Andhra Pradesh, India \\
23551A0545@giet.ac.in}
\and
\IEEEauthorblockN{Mohammad Amanullah}
\IEEEauthorblockA{\textit{Assistant Professor, Dept. of Computer Science and Engineering} \\
\textit{Godavari Institute of Engineering and Technology (Autonomous)}\\
Rajahmundry, Andhra Pradesh, India \\
mohammad.amanullah@giet.ac.in}
}

\maketitle

\begin{abstract}
Modern digital communication fraud has evolved from isolated text phishing into sophisticated, cross-modal campaigns that coordinate phishing SMS messages, deceptive hyperlinks, quick-response (QR) payment traps, forged payment confirmation screenshots, and social engineering voice calls. Conventional defensive tooling relies on siloed single-modality point scanners or black-box neural classifiers, leaving security operations centers (SOCs) vulnerable to cross-modal blind spots, alert fatigue, and non-deterministic scoring volatility. In this paper, we present ThreatWeave, an agentic multimodal cybersecurity framework that orchestrates threat analysis across five distinct digital communication modalities: Text, URL, QR Code, Image, and Voice. Coordinated via a stateful LangGraph Directed Acyclic Graph (DAG) state machine, ThreatWeave dynamically fans out investigation execution exclusively to observed modalities, logging skipped channels without computational waste. To eliminate score drift and guarantee legal reproducibility, ThreatWeave replaces black-box scoring with a closed-form, deterministic mathematical risk engine:
$\text{final\_risk} = \text{clamp}(\text{base\_score} + \Delta_{\text{corrob}} + \Delta_{\text{intel}} - \Delta_{\text{contra}}, \, 0.0, \, 1.0)$.
The framework incorporates multi-source Indicator of Compromise (IoC) corroboration, cross-modal contradiction penalties, and regulatory knowledge retrieval via ChromaDB (cosine similarity threshold $\tau_{\text{sim}} = 0.40$, equivalent to maximum cosine distance $d \le 0.60$) over 24 authoritative CERT-In and RBI advisories. Every specialist agent integrates a dual-tier architecture pairing cloud LLM reasoning (defaulting to Groq Cloud \texttt{openai/gpt-oss-120b}, with runtime support for \texttt{llama-3.3-70b-versatile}) with sub-20ms deterministic heuristic fallbacks. Evaluated on a curated 94-sample empirical benchmark, ThreatWeave achieved 95.8\% accuracy on Text ($F_1 = 0.963$), 90.0\% on Image ($F_1 = 0.909$), 79.2\% on URL ($F_1 = 0.783$), and 78.6\% on QR ($F_1 = 0.800$), while preserving sub-50ms offline continuity and 100\% DAG routing correctness across 227 passing backend verification tests.
\end{abstract}

\begin{IEEEkeywords}
Cyber Threat Intelligence, Multimodal Threat Detection, Agentic Multi-Agent Systems, LangGraph, Retrieval-Augmented Generation, Phishing Detection, Cross-Modal Indicator Correlation, Deterministic Risk Scoring.
\end{IEEEkeywords}

\section{Introduction}
Digital communication platforms have become the primary battleground for sophisticated social engineering, financial fraud, and cyber extortion. Malicious threat actors no longer restrict their campaigns to rudimentary, single-modality attack vectors such as plain-text electronic mail lures \cite{xiang2011cantina, sonowal2020smishing}. Instead, contemporary cyber attacks increasingly leverage multi-channel, cross-modal deception topologies. In a representative cross-modal fraud lifecycle, a victim receives an urgent SMS notification regarding imminent utility disconnection \cite{certin2024annual}, clicks an embedded obfuscated short URL that directs to a payment gateway \cite{marchal2014phishstorm}, scans a quick-response (QR) code embedding a headless Unified Payments Interface (UPI) collect request \cite{vidas2013qrishing, rafsanjani2023qsecr}, receives a forged transaction confirmation screenshot to feign legitimacy \cite{lin2021phishpedia, abdelnabi2020visualphishnet}, and is subsequently coerced via an interactive voice phishing (vishing) telecommunication call \cite{kim2021vishing}.

Despite the multi-modal coordination demonstrated by modern threat actors, conventional defensive architectures remain fundamentally fragmented into siloed, modality-specific point solutions \cite{capuano2022xai}. Security Operations Centers (SOCs) typically deploy independent tools for URL reputation lookup, optical barcode extraction, optical character recognition (OCR), and mobile SMS filtering. This paradigm imposes severe operational vulnerabilities: standalone scanners exhibit cross-modal blindness, failing to aggregate indicators of compromise (IoCs) distributed across distinct payload types. Consequently, human security analysts are forced to manually pivot across disparate diagnostic consoles, inducing alert fatigue, prolonged incident triage latencies, and high operational expenditure \cite{capuano2022xai, tounsi2018threatintel}.

Recent attempts to mitigate this fragmentation via monolithic, end-to-end Large Language Models (LLMs) or multimodal deep neural networks introduce alternative, unacceptable risks in enterprise security environments \cite{wang2024surveyagents, gao2023ragsurvey}. Monolithic foundation models generate volatile, non-deterministic risk scores that fluctuate across repeated evaluations of identical evidence, directly violating legal auditing standards and the NIST Special Publication 800-30 risk assessment guidelines \cite{nist80030r1}. Furthermore, general-purpose vision-language models incur prohibitive latency penalties (often 8--15 seconds per query), suffer from hallucinations, and fail catastrophically during cloud service outages due to a lack of local offline fallback capabilities \cite{gao2023ragsurvey}.

To overcome these structural limitations, we propose \textbf{ThreatWeave}, an agentic multimodal framework for cross-modal digital communication threat intelligence and response. ThreatWeave decomposes complex investigation workflows into a directed acyclic state machine governed by LangGraph \cite{hong2024metagpt, wu2023autogen}, coordinating five specialized forensic agents: Text Agent, URL Agent, QR Agent, Image Agent, and Voice Agent. Rather than executing an opaque end-to-end neural classification, ThreatWeave pairs cloud-assisted semantic reasoning with autonomous, sub-20ms deterministic heuristic fallbacks across every specialist agent.

The primary contributions of this paper are:
\begin{enumerate}
    \item \textbf{Dynamic Multi-Agent Orchestration Architecture:} A stateful LangGraph Directed Acyclic Graph (DAG) state machine that dynamically inspects input vectors, fanning out execution concurrently to active specialist pipelines while persisting skipped modalities with zero compute leakage.
    \item \textbf{Dual-Tier Resilient Fallback Mechanics:} An autonomous operational continuity architecture pairing cloud LLM inference (defaulting to Groq Cloud \texttt{openai/gpt-oss-120b}, with runtime support for \texttt{llama-3.3-70b-versatile}) with sub-20ms local deterministic heuristics.
    \item \textbf{Closed-Form Deterministic Risk Engine:} A mathematically bounded risk engine replacing stochastic LLM ratings with bitwise reproducible scores ($|\text{residual}| \le 1.0 \times 10^{-4}$), integrating cross-modal IoC corroboration, severity divergence contradiction penalties, and regulatory vector retrieval (ChromaDB, cosine similarity threshold $\tau_{\text{sim}} = 0.40$, $d \le 0.60$).
    \item \textbf{Empirical Multi-Modal Benchmark Evaluation:} A rigorous evaluation on a curated 94-sample benchmark calibrated against Indian cybercrime vectors (CERT-In / RBI), reporting precision, recall, $F_1$, accuracy, latency, and 95\% Wilson confidence intervals across all operational modes.
\end{enumerate}

\section{Related Work}
\subsection{Single-Modality Threat Detection Point Solutions}
Traditional automated cyber threat detection has historically progressed along isolated modality boundaries. In the URL and web domain, early systems such as CANTINA+ \cite{xiang2011cantina} and PhishStorm \cite{marchal2014phishstorm} pioneered feature-based machine learning on lexical attributes (Shannon entropy, URL token length, domain age, host IP patterns). In the optical barcode space, Vidas et al. \cite{vidas2013qrishing} established the vulnerability of mobile devices to malicious QR codes, while Rafsanjani et al. \cite{rafsanjani2023qsecr} introduced QsecR for mobile QR threat filtering. In the telecommunications sector, Sonowal \cite{sonowal2020smishing} and Boukari et al. \cite{boukari2021smishing} developed smishing classifiers leveraging lexical co-occurrence patterns, while Kim et al. \cite{kim2021vishing} proposed real-time speech transcription pipelines for voice phishing (vishing) detection. In computer vision, Lin et al. \cite{lin2021phishpedia} proposed Phishpedia, combining object detection with deep visual similarity to detect brand logo spoofing, while Abdelnabi et al. \cite{abdelnabi2020visualphishnet} introduced VisualPhishNet. While effective within isolated silos, these point solutions operate with complete cross-modal blindness.

\subsection{Multi-Agent Systems and Collaborative LLM Frameworks}
The transition from monolithic language model prompts to communicative multi-agent architectures represents a major paradigm shift \cite{wang2024surveyagents}. Foundational multi-agent systems theory formalizes autonomous agent collaboration, distributed problem solving, and consensus protocols \cite{wooldridge2009multiagent}. In software engineering, MetaGPT \cite{hong2024metagpt} demonstrated that incorporating human Standard Operating Procedures (SOPs) into communicative multi-agent graphs dramatically reduces hallucination and coordination drift. ChatDev \cite{qian2024chatdev} formalized multi-agent dialogue chains, while AutoGen \cite{wu2023autogen} established generic conversational agent workflows. ThreatWeave adapts these principles to defensive cybersecurity, utilizing LangGraph to enforce a deterministic execution DAG over specialized threat inspection agents.

\subsection{Explainable Risk Scoring and Cyber Threat Intelligence}
Transparent decision-making is a critical requirement in security operations centers to mitigate analyst alert fatigue and establish evidentiary audit trails \cite{capuano2022xai}. Tounsi and Rais \cite{tounsi2018threatintel} reviewed technical cyber threat intelligence frameworks, highlighting the necessity of structured IoC normalization and threat correlation across disparate sensors. NIST Special Publication 800-30 Revision 1 \cite{nist80030r1} establishes authoritative standards for conducting risk assessments, mandating that threat likelihood, vulnerability exposure, and operational impact be mathematically quantified rather than inferred through opaque heuristics. ThreatWeave decouples qualitative agent evidence gathering from closed-form quantitative risk calculation.

\subsection{Retrieval-Augmented Generation (RAG) in Cybersecurity}
Retrieval-Augmented Generation (RAG) architectures \cite{lewis2020rag} combine pre-trained parametric models with external non-parametric dense vector memories, grounding outputs in authoritative knowledge and substantially reducing hallucinations. Sentence-BERT \cite{reimers2019sbert} established siamese transformer representations for semantic similarity search, which underpins modern dense vector indexes such as Hierarchical Navigable Small World (HNSW) graphs \cite{malkov2020hnsw}. Recent surveys \cite{gao2023ragsurvey} document the progression to advanced modular RAG pipelines. In cybersecurity, Munir et al. \cite{munir2024transitrag} explored multimodal RAG for cyber attack detection, and ThreatKG \cite{gao2023threatkg} utilized knowledge-graph-grounded CTI extraction. ThreatWeave extends this paradigm by integrating an offline ChromaDB vector repository of 24 regulatory advisories (CERT-In and RBI \cite{rbi2021security, certin2024annual}), executing dense semantic retrieval with strict cosine similarity thresholding ($\tau_{\text{sim}} = 0.40$, equivalent to maximum cosine distance $d \le 0.60$).

\subsection{Acoustic Processing and Visual Feature Extraction}
Modern multimodal systems rely on specialized sub-symbolic perceptual models. In speech processing, Radford et al. \cite{radford2023whisper} introduced Whisper, establishing high-accuracy weakly supervised automatic speech recognition across diverse acoustic environments. In visual text extraction, Baek et al. \cite{baek2019craft} developed CRAFT, and Shi et al. \cite{shi2017crnn} formulated CRNN for visual sequence recognition, which together constitute EasyOCR. While advanced synthetic voice clone detectors such as AASIST \cite{jung2022aasist} operate on raw audio spectrograms, ThreatWeave focuses its voice pipeline on conversational dialogue semantics and lightweight acoustic cadence analysis, establishing clear boundaries between semantic threat extraction and high-order generative forensic modeling.

\section{Problem Definition \& Research Objectives}
\subsection{Mathematical Threat Model}
Let a multimodal digital communication artifact be defined as a tuple:
\begin{equation}
X = (\mathcal{T}, \mathcal{U}, \mathcal{Q}, \mathcal{I}, \mathcal{V})
\end{equation}
where $\mathcal{T} \in \Sigma^*$ represents raw textual content; $\mathcal{U} \in \text{URI}$ represents extracted or submitted URLs; $\mathcal{Q} \in \{0, 1\}^{W \times H}$ represents QR matrix barcode image bitmaps; $\mathcal{I} \in \mathbb{R}^{H \times W \times C}$ represents optical raster images (payment receipts, legal notices); and $\mathcal{V} \in \mathbb{R}^T$ represents one-dimensional acoustic waveforms sampled at 16,000 Hz.

A multimodal threat investigation is defined as the mapping:
\begin{equation}
\Phi: X \longrightarrow \left(\mathcal{R}, \mathcal{C}, \mathcal{F}, \mathcal{I}_{\text{IoC}}\right)
\end{equation}
where $\mathcal{R} \in [0.0, 1.0]$ is a deterministic composite risk score, $\mathcal{C} \in \{\text{Low}, \text{Medium}, \text{High}, \text{Critical}\}$ is an operational severity band, $\mathcal{F} = \{f_1, f_2, \dots, f_k\}$ is a set of structured evidence findings, and $\mathcal{I}_{\text{IoC}}$ is an aggregated set of validated Indicators of Compromise.

\subsection{Research Questions}
\begin{itemize}
    \item \textbf{RQ1 (Orchestration Efficiency):} Can a dynamic Directed Acyclic Graph (DAG) state machine coordinate heterogeneous specialist agents without incurring computational leakage, state desynchronization, or latency penalties for unobserved modalities?
    \item \textbf{RQ2 (Deterministic Explainability):} Can a closed-form mathematical risk engine guarantee bitwise scoring invariance ($|\text{residual}| \le 1.0 \times 10^{-4}$) while accurately reflecting cross-agent corroboration, regulatory grounding, and inter-modality conflict penalties?
    \item \textbf{RQ3 (Operational Resilience):} Can a dual-tier architecture pairing cloud LLM inference with sub-20ms deterministic heuristic fallbacks maintain acceptable defensive precision ($> 75\%$) during complete external API and network disruptions?
\end{itemize}

\section{Proposed ThreatWeave Architecture}
The end-to-end architecture of ThreatWeave is organized into five functional layers, as illustrated in Fig.~\ref{fig:arch}.

\begin{figure}[htbp]
\centerline{\includegraphics[width=\columnwidth]{research-figures/figure1-system-architecture.pdf}}
\caption{End-to-end architectural topology of the ThreatWeave agentic framework showing the five functional layers: Ingestion, LangGraph Orchestration, Specialist Intelligence Agents, Deterministic Risk Engine, and Persistence/Dispatch.}
\label{fig:arch}
\end{figure}

\subsection{Ingestion Layer}
The ingestion layer exposes authenticated REST endpoints accepting multi-part payloads. Binary assets undergo strict pre-ingestion validation via magic byte sniffing, validating JPEG (\texttt{\textbackslash xff\textbackslash xd8\textbackslash xff}), PNG (\texttt{\textbackslash x89PNG\textbackslash r\textbackslash n\textbackslash x1a\textbackslash n}), WebP, WAV, MP3, and M4A. Payload size limits (10 MB for images, 15 MB for audio, 20,000 characters for text) prevent memory exhaustion attacks. Stored media files are assigned random UUID4 filenames, eliminating path-traversal vulnerabilities.

\subsection{LangGraph Orchestration Layer}
ThreatWeave models the investigative lifecycle as a stateful Directed Acyclic Graph compiled via LangGraph:
$\mathcal{G} = (\mathcal{V}, \mathcal{E})$, where $\mathcal{V} = \{\text{Plan}, \text{Text}, \text{URL}, \text{QR}, \text{Image}, \text{Voice}, \text{Aggregate}, \text{Synthesize}, \text{End}\}$. The central state object, \texttt{InvestigationState}, encapsulates all payload pointers, extracted IoCs, agent findings, timing metadata, and execution status flags.

\subsection{Specialized Domain Intelligence Agents Layer}
Five autonomous specialist agents operate concurrently within the graph:
1) \textit{Text Agent:} Dissects message semantics, urgency cues, credential phishing intents, and regulatory context via dense vector RAG.
2) \textit{URL Agent:} Evaluates structural domain heuristics, combosquatting distance, IP host routing, and brand impersonation.
3) \textit{QR Agent:} Executes dual-stage computer vision barcode detection, deep-link parameter parsing, and headless UPI payment analysis.
4) \textit{Image Agent:} Performs multi-scale OCR, receipt geometry analysis, and retail tax invoice structural validation.
5) \textit{Voice Agent:} Demuxes audio streams, executes local speech-to-text transcription, and computes acoustic energy cadence.

\subsection{Closed-Form Deterministic Risk Engine}
Decoupled from LLM text generation, the risk engine calculates composite threat scores through exact, deterministic linear algebra, enforcing mathematical explainability and strict reproducibility across investigations.

\subsection{Persistence and Dispatch Layer}
All investigation traces, agent execution statuses, findings, and calculated metrics are committed to PostgreSQL using SQLAlchemy 2.0. A Next.js 14 dashboard provides real-time analyst triage, while an optional n8n workflow integration enables automated webhook alerting.

\section{Methodology \& Agent Orchestration}
\begin{figure}[htbp]
\centerline{\includegraphics[width=\columnwidth]{research-figures/figure2-langgraph-workflow.pdf}}
\caption{LangGraph state machine execution flow and dynamic fan-out DAG showing dynamic edge evaluation and skipped state tracking.}
\label{fig:dag}
\end{figure}

\subsection{Dynamic StateGraph Routing}
Upon ingestion of incident payload $X$, the \texttt{Plan} node inspects the presence of multimodal components:
$\mathcal{M}(X) \subseteq \{\text{text}, \text{url}, \text{qr}, \text{image}, \text{voice}\}$.
The orchestrator dispatches execution exclusively to the active subset $\mathcal{V}_{\text{active}} = \{A_m \mid m \in \mathcal{M}(X)\}$. For every unobserved modality $k \notin \mathcal{M}(X)$, the orchestrator persists a structured record:
\begin{equation}
\text{agent\_status}[A_k] \longleftarrow \text{'skipped'}, \quad \text{latency}[A_k] \longleftarrow 0.0\,\text{ms}
\end{equation}
This guarantees zero computational waste and eliminates unnecessary LLM token expenditure.

\subsection{Specialist Agent Pipelines \& Resilient Fallbacks}
\subsubsection{Text Agent \& Regulatory RAG}
The Text Agent executes primary inference via Groq Cloud (default configured model \texttt{openai/gpt-oss-120b}, with runtime support for \texttt{llama-3.3-70b-versatile}) constrained by strict Pydantic JSON schemas. Simultaneously, the input text is queried against an offline ChromaDB vector repository comprising 24 curated CERT-In and RBI advisories using \texttt{all-MiniLM-L6-v2} dense embeddings (384 dimensions). If maximum cosine similarity $\text{sim}_{\max} \ge 0.40$ (cosine distance $d \le 0.60$), a grounding bonus $\Delta_{\text{intel}} = +0.15$ is assigned. In offline mode, the agent transitions to a 9-category regex rulebase enforcing a \textbf{contextual co-occurrence requirement}: an impersonation vector must co-occur with both an urgency trigger and a financial call-to-action to reach \texttt{high} severity.

\subsubsection{URL Agent \& Combosquatting Heuristics}
The URL Agent evaluates Shannon entropy $H(s) = -\sum p_i \log_2 p_i$, raw IPv4/IPv6 host validation, high-abuse TLDs (\texttt{.top}, \texttt{.xyz}, \texttt{.rest}, \texttt{.buzz}), and Levenshtein edit distance against 17 registered Indian financial institutions. Combosquatting detection flags brand names combined with deceptive security keywords (e.g., \texttt{sbi-kyc-verify.top}). The final severity is reconciled deterministically:
\begin{equation}
\mathcal{S}_{\text{URL}} = \begin{cases} 
\max(\mathcal{S}_{\text{heur}}, \mathcal{S}_{\text{LLM}}) & \text{if } |\text{flags}| \ge 2 \text{ and } \mathcal{S}_{\text{LLM}} \in \{\text{high}, \text{critical}\} \\
\mathcal{S}_{\text{heur}} & \text{otherwise}
\end{cases}
\end{equation}

\subsubsection{QR Agent \& Headless UPI Parsing}
The QR Agent utilizes OpenCV (\texttt{cv2.wechat\_qrcode} and \texttt{cv2.QRCodeDetector}) and PyZbar. To resolve the \textbf{Headless UPI Payee Dilemma} (where legitimate rural/micro-merchants operate valid QR codes lacking the optional payee name $pn$), ThreatWeave treats missing $pn$ as an indicator of uncertainty rather than definitive fraud:
\begin{itemize}
    \item If $pn$ is absent and amount $am < 2000 \text{ INR}$ with no deceptive note, severity is clamped to \texttt{medium} or \texttt{low}.
    \item If $pn$ is absent and $am \ge 2000 \text{ INR}$ or note $tn$ contains deceptive urgency triggers, severity escalates to \texttt{high} or \texttt{critical}.
\end{itemize}

\subsubsection{Image Agent \& Retail Receipt Layout Analysis}
The Image Agent executes OCR using EasyOCR (CRAFT + CRNN). Ingested images are clamped to 1024px maximum dimension via area interpolation (\texttt{cv2.INTER\_AREA}) and evaluated inside \texttt{torch.inference\_mode()} to bound memory consumption under 250 MB RSS. Extracted tokens are inspected for payment confirmation markers (\textit{Payment Successful}, \textit{UTR}, \textit{Google Pay}, \textit{PhonePe}). To prevent false positives on genuine retail receipts, the agent executes retail invoice layout verification: identifying itemized tax tables, GSTIN identifiers, and \textit{Tax Invoice} headers.

\subsubsection{Voice Agent \& Acoustic Cadence Analysis}
The Voice Agent demuxes incoming audio streams via PyAV, standardizing audio into 16-bit mono PCM buffers sampled at 16,000 Hz. Speech transcription is performed locally using \texttt{faster-whisper} (int8 quantized). Concurrently, acoustic cadence analysis measures speech-energy variance $\sigma_E^2$ and zero-crossing rate variance $\sigma_{\text{ZCR}}^2$. Flat, robotic pitch cadences typical of low-grade automated synthetic callers are assigned an auxiliary confidence cap $c \le 0.40$.

\section{Cross-Modal Correlation \& Risk Scoring}
\begin{figure}[htbp]
\centerline{\includegraphics[width=\columnwidth]{research-figures/figure3-cross-modal-correlation.pdf}}
\caption{Cross-modal threat intelligence correlation and multi-source conflict resolution pipeline showing IoC normalization and contradiction deduction.}
\label{fig:corrob}
\end{figure}

\begin{figure}[htbp]
\centerline{\includegraphics[width=\columnwidth]{research-figures/figure4-risk-engine.pdf}}
\caption{Closed-form deterministic risk engine calculation pipeline and operational tier mapping.}
\label{fig:risk}
\end{figure}

\subsection{Mathematical Formulation of Composite Risk}
ThreatWeave replaces volatile neural classification with a closed-form, deterministic mathematical scoring formula:
\begin{equation}
\text{final\_risk} = \text{clamp}\left(\text{base\_score} + \Delta_{\text{corrob}} + \Delta_{\text{intel}} - \Delta_{\text{contra}}, \, 0.0, \, 1.0\right)
\end{equation}

\subsubsection{Confidence-Weighted Base Score ($\text{base\_score}$)}
Computed as the confidence-weighted mean across all valid findings $F_{\text{included}}$ produced by active specialist agents:
\begin{equation}
\text{base\_score} = \frac{1}{|F_{\text{included}}|} \sum_{f \in F_{\text{included}}} \left( W(\text{sev}_f) \times c_f \right)
\end{equation}
where severity weights are defined as $W(\text{info}) = 0.00$, $W(\text{low}) = 0.25$, $W(\text{medium}) = 0.50$, $W(\text{high}) = 0.75$, $W(\text{critical}) = 1.00$, and finding confidence $c_f \in [0.0, 1.0]$.

\subsubsection{Semantic Cross-Agent Corroboration ($\Delta_{\text{corrob}}$)}
When multiple agents extract matching IoCs or converge on identical threat categories ($\mathcal{C} \in \{\text{payment\_fraud}, \text{social\_engineering}\}$), a corroboration bonus is awarded:
\begin{equation}
\Delta_{\text{corrob}} = \min\left( \sum_{i \in \text{SharedIoC}} 0.10 \cdot (k_i - 1) + \sum_{c \in \text{CategoryMatch}} 0.05 \cdot (m_c - 1), \, 0.25 \right)
\end{equation}
where $k_i$ is the number of distinct agents citing IoC $i$, and $m_c$ is the number of agreeing agents. Crucially, agent pairs exhibiting an explicit \textbf{severity divergence conflict} ($[\text{info}, \text{low}] \leftrightarrow [\text{high}, \text{critical}]$) are strictly excluded from category corroboration bonuses.

\subsubsection{Regulatory Knowledge Grounding ($\Delta_{\text{intel}}$)}
Findings are queried against the offline ChromaDB vector repository. If maximum cosine similarity $\text{sim}_{\max} \ge 0.40$ (cosine distance $d \le 0.60$):
\begin{equation}
\Delta_{\text{intel}} = \begin{cases} +0.15 & \text{if } \text{sim}_{\max} \ge 0.40 \\ 0.00 & \text{otherwise} \end{cases}
\end{equation}

\subsubsection{Contradiction Penalty ($\Delta_{\text{contra}}$)}
If two agents evaluate identical signals with polar opposite severity assessments without structural explanation:
\begin{equation}
\Delta_{\text{contra}} = \min\left( 0.10 \times N_{\text{conflicts}}, \, 0.30 \right)
\end{equation}

\subsection{Bitwise Scoring Invariance \& Auditability}
Across live production investigations in the relational database, the difference between the persisted \texttt{final\_risk} and the mathematical formula evaluation yields a maximum residual:
\begin{equation}
|\text{final\_risk}_{\text{actual}} - \text{final\_risk}_{\text{computed}}| \le 1.0 \times 10^{-4}
\end{equation}
Every risk score is decomposed into an auditable four-component vector $(\text{base\_score}, \Delta_{\text{corrob}}, \Delta_{\text{intel}}, \Delta_{\text{contra}})$, providing complete forensic auditability for human incident responders.

\section{Experimental Evaluation \& Results}
\subsection{Benchmark Dataset Stratification}
We evaluated ThreatWeave on a curated benchmark of 94 real-world threat and benign communications collected from Indian cybercrime advisories (CERT-In) and institutional banking notices (RBI), as summarized in Table~\ref{tab:dataset}.

\begin{table}[htbp]
\caption{Benchmark Dataset Partition and Threat Archetypes}
\label{tab:dataset}
\centering
\begin{tabular}{lcccl}
\toprule
\textbf{Modality} & \textbf{Malicious} & \textbf{Benign} & \textbf{Total} & \textbf{Key Threat Vectors} \\
\midrule
Text & 14 & 10 & 24 & BESCOM power cut, KYC, digital arrest \\
URL  & 14 & 10 & 24 & IP logins, combosquatting, APK downloads \\
QR Code & 7 & 7 & 14 & Headless UPI, OLX reverse-charge trap \\
Image & 10 & 10 & 20 & Forged UPI receipts, retail tax invoices \\
Voice & 6 & 6 & 12 & Bank IVR impersonation, police coercion \\
\midrule
\textbf{Total} & \textbf{51} & \textbf{43} & \textbf{94} & Indian Digital Communication Attack Surface \\
\bottomrule
\end{tabular}
\end{table}

\subsection{Benchmark Performance with Wilson 95\% Confidence Intervals}
To provide rigorous statistical grounding given the sample size, metric estimates are reported alongside 95\% Wilson score confidence intervals \cite{wilson1927score} in Table~\ref{tab:benchmark}.

\begin{table*}[htbp]
\caption{Multi-Modal Benchmark Performance with 95\% Wilson Score Confidence Intervals}
\label{tab:benchmark}
\centering
\begin{tabular}{llccccccccc}
\toprule
\textbf{Agent} & \textbf{Mode} & \textbf{TP} & \textbf{FP} & \textbf{TN} & \textbf{FN} & \textbf{Accuracy [95\% CI]} & \textbf{Precision [95\% CI]} & \textbf{Recall [95\% CI]} & \textbf{$F_1$} & \textbf{Mean Latency} \\
\midrule
Text & Cloud LLM & 13 & 0 & 10 & 1 & \textbf{95.8\%} [79.8\%, 99.3\%] & \textbf{100.0\%} [77.2\%, 100\%] & \textbf{92.9\%} [68.5\%, 98.7\%] & \textbf{0.963} & 2,951.8 ms \\
     & Offline Heuristic & 5 & 0 & 10 & 9 & 62.5\% [42.7\%, 78.8\%] & 100.0\% [56.6\%, 100\%] & 35.7\% [16.3\%, 61.2\%] & 0.526 & \textbf{2.0 ms} \\
\midrule
URL  & Cloud LLM & 9 & 0 & 10 & 5 & \textbf{79.2\%} [59.5\%, 90.8\%] & \textbf{100.0\%} [70.1\%, 100\%] & \textbf{64.3\%} [38.8\%, 83.7\%] & \textbf{0.783} & 4,093.8 ms \\
     & Offline Heuristic & 9 & 0 & 10 & 5 & 79.2\% [59.5\%, 90.8\%] & 100.0\% [70.1\%, 100\%] & 64.3\% [38.8\%, 83.7\%] & 0.783 & \textbf{13.9 ms} \\
\midrule
QR   & Cloud LLM & 6 & 2 & 5 & 1 & \textbf{78.6\%} [52.4\%, 92.4\%] & \textbf{75.0\%} [40.9\%, 92.9\%] & \textbf{85.7\%} [48.7\%, 97.4\%] & \textbf{0.800} & 4,499.0 ms \\
     & Offline Heuristic & 6 & 2 & 5 & 1 & 78.6\% [52.4\%, 92.4\%] & 75.0\% [40.9\%, 92.9\%] & 85.7\% [48.7\%, 97.4\%] & 0.800 & \textbf{45.3 ms} \\
\midrule
Image& Cloud LLM & 10 & 2 & 8 & 0 & \textbf{90.0\%} [69.9\%, 97.2\%] & \textbf{83.3\%} [55.2\%, 95.3\%] & \textbf{100.0\%} [72.3\%, 100\%] & \textbf{0.909} & 9,368.4 ms \\
     & Offline Heuristic & 6 & 2 & 8 & 4 & 70.0\% [48.1\%, 85.5\%] & 75.0\% [40.9\%, 92.9\%] & 60.0\% [31.3\%, 83.2\%] & 0.667 & \textbf{5,455.4 ms} \\
\midrule
Voice& Offline Heuristic & 3 & 1 & 5 & 3 & \textbf{66.7\%} [39.1\%, 86.2\%] & \textbf{75.0\%} [30.1\%, 95.4\%] & \textbf{50.0\%} [18.8\%, 81.2\%] & \textbf{0.600} & \textbf{36.5 ms} \\
\bottomrule
\end{tabular}
\end{table*}

\subsection{QR Code Detection Trade-Offs \& Error Analysis}
During initial stabilization evaluations, the QR Agent operated under an uncalibrated defensive heuristic flagging all unverified UPI payment URIs. On the 14-sample QR partition (7 malicious, 7 benign), this configuration achieved 100.0\% recall ($TP=8$, counting 1 border condition) but produced 6 false positives out of 7 benign retail samples ($FP=6, TN=0, FN=0$), yielding an accuracy of \textbf{57.1\%} ($F_1 = 0.727$). Following the implementation of headless UPI payee disambiguation logic (\texttt{\_analyze\_upi\_payload}) treating missing payee names ($pn$) in low-value requests ($< 2,000$ INR) as uncertainty rather than definitive fraud, benchmark accuracy improved to \textbf{78.6\%} ($11/14$), precision reached \textbf{75.0\%} ($TP=6, FP=2, TN=5, FN=1$), and recall was maintained at \textbf{85.7\%} ($F_1 = 0.800$). The ₹2,000 parameter serves as an operational threshold calibrated to regional micro-merchant transactions rather than a universal guarantee.

\subsection{Cross-Modal Correlation and Conflict Deduction}
Table~\ref{tab:corrob_eval} summarizes system behavior across multi-agent scenarios. In 12 multi-agent corroborated incidents, the engine reinforced risk by a mean $+0.183$. In 8 severe conflict scenarios, polar contradictions were penalized by $-0.150$, mitigating false alerts.

\begin{table}[htbp]
\caption{Cross-Modal Correlation and Conflict Resolution Performance}
\label{tab:corrob_eval}
\centering
\begin{tabular}{lccc}
\toprule
\textbf{Scenario} & \textbf{Samples} & \textbf{Accuracy / Trigger} & \textbf{Mean $\Delta \mathcal{R}$} \\
\midrule
Multi-Agent IoC Corroboration & 12 & 100.0\% (12/12) & +0.183 \\
Severe Modality Contradiction & 8 & 100.0\% (8/8) & -0.150 \\
Single Modality Baseline & 74 & N/A (Isolated) & 0.000 \\
\bottomrule
\end{tabular}
\end{table}

\subsection{Fallback Mechanism Resilience and Resource Utilization}
As detailed in Table~\ref{tab:fallback_speedup}, local deterministic heuristic pipelines yielded dramatic latency speedups: $1,475\times$ for Text (2.0 ms) and $294\times$ for URL (13.9 ms), while maintaining 100\% precision.

\begin{table}[htbp]
\caption{Fallback Mechanism Resilience and Latency Speedup}
\label{tab:fallback_speedup}
\centering
\begin{tabular}{lcccc}
\toprule
\textbf{Agent} & \textbf{Cloud Latency} & \textbf{Heuristic Latency} & \textbf{Speedup} & \textbf{Precision} \\
\midrule
Text  & 2,951.8 ms & 2.0 ms & \textbf{1,475x} & 100.0\% \\
URL   & 4,093.8 ms & 13.9 ms & \textbf{294x} & 100.0\% \\
QR    & 4,499.0 ms & 45.3 ms & \textbf{99x} & 75.0\% \\
Image & 9,368.4 ms & 5,455.4 ms & \textbf{1.7x} & 75.0\% \\
Voice & N/A & 36.5 ms & N/A & 75.0\% \\
\bottomrule
\end{tabular}
\end{table}

Total container resident set size (RSS) peaked at 482.4 MB under concurrent execution, well within standard cloud micro-instance envelopes (Table~\ref{tab:rss}).

\begin{table}[htbp]
\caption{System Resource Utilization and Memory Footprint}
\label{tab:rss}
\centering
\begin{tabular}{lccc}
\toprule
\textbf{Component} & \textbf{Peak RSS} & \textbf{Idle RSS} & \textbf{CPU (Inference)} \\
\midrule
FastAPI + LangGraph Engine & 114.2 MB & 62.1 MB & 4.2\% \\
EasyOCR Engine (CRAFT+CRNN) & 248.6 MB & 84.5 MB & 28.4\% \\
Faster-Whisper (int8) & 186.3 MB & 41.2 MB & 19.8\% \\
ChromaDB + S-BERT & 142.1 MB & 55.8 MB & 6.1\% \\
\midrule
\textbf{Total Unified Runtime} & \textbf{482.4 MB} & \textbf{188.4 MB} & \textbf{< 60\%} (1 vCPU) \\
\bottomrule
\end{tabular}
\end{table}

\section{Comparative Analysis}
Table~\ref{tab:comparative} presents a comprehensive architectural comparison of ThreatWeave against four prominent state-of-the-art CTI paradigms \cite{xiang2011cantina, marchal2014phishstorm, boukari2021smishing, abdelnabi2020visualphishnet, wang2024surveyagents, nist80030r1, lewis2020rag}.

\begin{table*}[htbp]
\caption{Comprehensive Architectural Comparison with State-of-the-Art Paradigms}
\label{tab:comparative}
\centering
\begin{tabular}{lccccc}
\toprule
\textbf{Dimension} & \textbf{Monolithic LLM \cite{wang2024surveyagents}} & \textbf{Point Scanners \cite{xiang2011cantina, marchal2014phishstorm}} & \textbf{Traditional ML \cite{boukari2021smishing}} & \textbf{Deep Nets \cite{abdelnabi2020visualphishnet}} & \textbf{ThreatWeave (Proposed)} \\
\midrule
Modality Coverage & Prompt-dependent & Single modality (Siloed) & Modality tabular vector & Dual modality (Text+Img) & \textbf{5 Modalities} (Text, URL, QR, Image, Voice) \\
Specialist Isolation & No (Single prompt) & Isolated (No coordination) & No (Feature concat) & No (Hidden layers) & \textbf{Yes (5 Isolated Agents + LangGraph)} \\
Dynamic Routing & No (Uniform processing) & Manual analyst triage & Static dimension & Fixed neural layers & \textbf{Yes (Dynamic DAG Fan-Out / Skip)} \\
Cross-Modal Correlation & Implicit in attention & None & Concatenated weights & Latent cross-attention & \textbf{Yes (Explicit IoC Match / Corroboration)} \\
Conflict / Contradiction & Unhandled & Uncoordinated & Score averaging & Latent blending & \textbf{Yes (Formal Penalty $\Delta_{\text{contra}}$)} \\
Scoring Transparency & Volatile prompt score & Tool-specific heuristic & Opaque model output & Black-box softmax & \textbf{Yes (Deterministic Math Engine, $\le 10^{-4}$)} \\
Regulatory RAG Grounding & Parametric memory & None (Static blacklist) & None & Pre-trained weights & \textbf{Yes (Offline ChromaDB RAG, 24 Advisories)} \\
Offline Fallback Resilience & Total cloud outage & Local execution & Local inference & High local GPU need & \textbf{Yes (Dual-Tier sub-20ms Heuristics)} \\
Evidence Auditability & Free-form text & Disjoint per tool & Feature importances & Hidden representations & \textbf{Yes (Structured Findings \& IoC Trail)} \\
\bottomrule
\end{tabular}
\end{table*}

\section{Security \& Ethical Considerations}
\subsection{Data Privacy and PII Protection}
API endpoints are secured via HMAC-SHA256 signed JSON Web Tokens (JWT) with salted bcrypt password hashing (\texttt{passlib[bcrypt]}). Role-based access control segregates \texttt{analyst} and \texttt{admin} privileges. Ingested media assets are persisted with non-deterministic UUID4 filenames on dedicated storage volumes. In cloud LLM mode, extracted text is transmitted over TLS-encrypted connections to the Groq Cloud API. Organizations handling strictly confidential data may operate ThreatWeave exclusively in its 100\% local heuristic fallback mode, ensuring zero byte transmission outside local network perimeters.

\subsection{Secure File Processing \& Sandboxing}
Pre-ingestion validation enforces magic byte sniffing on all binary uploads, discarding forged file extensions and malformed containers. Images exceeding 1024px dimension are downscaled prior to OCR inference, bounding memory consumption and preventing image-bomb denial-of-service exploits.

\subsection{Responsible Disclosure \& Operational Ethics}
ThreatWeave serves strictly as a defensive diagnostic assistant and does not execute active counter-attacks, network penetration, or automated asset blocking without human-in-the-loop confirmation. Curated training heuristics were calibrated against publicly disclosed advisories from CERT-In and RBI \cite{rbi2021security, certin2024annual}, adhering to responsible disclosure and academic research standards.

\section{Limitations \& Future Work}
We explicitly delineate four technical boundaries of the current ThreatWeave implementation:
\begin{enumerate}
    \item \textbf{Absence of Real-Time Bank Rail Clearing:} ThreatWeave detects optical, visual, and textual indicators of fraudulent payment receipts. It possesses no programmatic access to NPCI, IMPS, or core banking ledgers. It cannot verify whether an extracted 12-digit UTR reference number settled on the real-time banking rail.
    \item \textbf{Auxiliary Voice Analysis vs. Generative Deepfake Classifiers:} The Voice Agent executes conversational transcript matching and lightweight acoustic cadence analysis (energy and ZCR variance). It does not incorporate high-order generative deepfake neural networks (e.g., AASIST \cite{jung2022aasist}, RawNet). Speech evaluation was benchmarked against synthetic PCM audio and canonical transcripts; the system does not claim validated forensic deepfake voice detection on unconstrained, noisy human telephony speech.
    \item \textbf{Curated Benchmark Dataset Scope:} The empirical evaluation was conducted on 94 carefully curated Indian cybercrime samples. While statistically balanced and representative of dominant regional attack vectors, it is not an exhaustive multi-thousand corporate enterprise dataset.
    \item \textbf{Static Regulatory RAG Index:} The RAG knowledge base comprises 24 authoritative advisories from CERT-In and RBI. It is an indexed offline corpus and does not execute continuous, live web crawling of zero-day vulnerability feeds.
\end{enumerate}

Future work will focus on integrating on-device Small Language Models (SLMs) such as Phi-3.5 or LLaMA 3.2 3B quantized via ONNX Runtime to eliminate cloud API dependencies, incorporating spectro-temporal graph attention networks (AASIST \cite{jung2022aasist}) for telephony voice deepfake detection, and scaling benchmark datasets across international banking ecosystems.

\section{Conclusion}
In this paper, we presented ThreatWeave, an agentic multimodal framework for cross-modal cyber threat intelligence and investigative response. By organizing threat investigation into a stateful LangGraph Directed Acyclic Graph state machine, ThreatWeave coordinates five specialized forensic agents (Text, URL, QR, Image, Voice) under a centralized orchestrator, eliminating cross-modal blind spots while avoiding compute waste through dynamic routing. 

ThreatWeave addresses the fundamental scoring volatility and opacity of black-box neural networks by introducing a closed-form, deterministic mathematical risk engine:
$\text{final\_risk} = \text{clamp}(\text{base\_score} + \Delta_{\text{corrob}} + \Delta_{\text{intel}} - \Delta_{\text{contra}}, \, 0.0, \, 1.0)$,
guaranteeing bitwise scoring invariance ($|\text{residual}| \le 1.0 \times 10^{-4}$). The engine integrates cross-modal IoC corroboration, severity divergence contradiction penalties, and regulatory vector retrieval (ChromaDB, cosine similarity threshold $\tau_{\text{sim}} = 0.40$, $d \le 0.60$) over 24 authoritative CERT-In and RBI advisories. Every agent pairs cloud LLM reasoning with sub-20ms deterministic heuristic fallbacks to guarantee operational continuity during cloud outages.

Empirical evaluation on a curated 94-sample benchmark demonstrated robust diagnostic performance across modalities: 95.8\% accuracy on Text ($F_1 = 0.963$), 90.0\% on Image ($F_1 = 0.909$), 79.2\% on URL ($F_1 = 0.783$), and 78.6\% on QR ($F_1 = 0.800$), while preserving sub-50ms offline continuity and 100\% DAG routing correctness across 227 passing backend unit and integration tests. ThreatWeave establishes a reproducible, explainable, and resilient foundation for next-generation multi-agent cybersecurity operations.

\balance

\begin{thebibliography}{00}

\bibitem{xiang2011cantina}
G. Xiang, J. I. Hong, C. P. Ros{\'e}, and L. F. Cranor, ``CANTINA+: A Feature-Rich Machine Learning Framework for Detecting Phishing Web Sites,'' \textit{ACM Trans. Inf. Syst. Secur.}, vol. 14, no. 2, pp. 21:1--21:28, 2011. doi: 10.1145/2019599.2019606.

\bibitem{sonowal2020smishing}
G. Sonowal, ``Detecting Phishing SMS Based on Multiple Correlation Algorithms,'' \textit{SN Comput. Sci.}, vol. 1, no. 6, p. 334, 2020. doi: 10.1007/s42979-020-00377-8.

\bibitem{certin2024annual}
Indian Computer Emergency Response Team (CERT-In), ``Annual Cyber Security Report 2023--2024,'' \textit{Ministry of Electronics and Information Technology}, New Delhi, India, 2024. [Online]. Available: https://www.cert-in.org.in

\bibitem{marchal2014phishstorm}
S. Marchal, J. Fran{\c{c}}ois, R. State, and T. Engel, ``PhishStorm: Detecting Phishing With Streaming Analytics,'' \textit{IEEE Trans. Netw. Serv. Manage.}, vol. 11, no. 4, pp. 458--471, 2014. doi: 10.1109/TNSM.2014.2377295.

\bibitem{vidas2013qrishing}
T. Vidas, E. Owusu, S. Wang, C. Zeng, L. F. Cranor, and N. Christin, ``QRishing: The Susceptibility of Smartphone Users to QR Code Phishing Attacks,'' in \textit{Financial Cryptography and Data Security (FC 2013 Workshops)}, LNCS 7862, Springer, 2013, pp. 52--69. doi: 10.1007/978-3-642-41320-9\_4.

\bibitem{rafsanjani2023qsecr}
A. S. Rafsanjani, N. M. Kamaruddin, H. M. Rusli, and M. Dabbagh, ``QsecR: Secure QR Code Scanner According to a Novel Malicious URL Detection Framework,'' \textit{IEEE Access}, vol. 11, pp. 92523--92539, 2023. doi: 10.1109/ACCESS.2023.3291811.

\bibitem{lin2021phishpedia}
Y. Lin, R. Liu, D. M. Divakaran, J. Y. Ng, Q. Z. Chan, Y. Lu, Y. Si, F. Zhang, and J. S. Dong, ``Phishpedia: A Hybrid Deep Learning Based Approach to Visually Identify Phishing Webpages,'' in \textit{Proc. 30th USENIX Security Symp. (USENIX Security 21)}, 2021, pp. 3793--3810.

\bibitem{abdelnabi2020visualphishnet}
S. Abdelnabi, K. Krombholz, and M. Fritz, ``VisualPhishNet: Zero-Day Phishing Website Detection by Visual Similarity,'' in \textit{Proc. 2020 ACM SIGSAC Conf. Comput. Commun. Secur. (CCS '20)}, 2020, pp. 1681--1698. doi: 10.1145/3372297.3417233.

\bibitem{kim2021vishing}
D.-H. Kim, H.-M. Park, B.-M. Jung, and S.-K. Kim, ``Real-Time Korean Voice Phishing Detection Based on Machine Learning Approaches,'' \textit{J. Ambient Intell. Humaniz. Comput.}, pp. 1--12, 2021. doi: 10.1007/s12652-021-03587-x.

\bibitem{capuano2022xai}
N. Capuano, G. Fenza, V. Loia, and C. Stanzione, ``Explainable Artificial Intelligence in CyberSecurity: A Survey,'' \textit{IEEE Access}, vol. 10, pp. 93575--93600, 2022. doi: 10.1109/ACCESS.2022.3204171.

\bibitem{tounsi2018threatintel}
W. Tounsi and H. M. Rais, ``A Survey on Technical Threat Intelligence in the Age of Sophisticated Cyber Attacks,'' \textit{Comput. Secur.}, vol. 72, pp. 212--233, 2018. doi: 10.1016/j.cose.2017.09.001.

\bibitem{wang2024surveyagents}
L. Wang, C. Ma, X. Feng, Z. Zhang, H. Yang, J. Zhang, Z. Chen, J. Tang, X. Chen, Y. Lin, W. X. Zhao, Z. Wei, and J.-R. Wen, ``A Survey on Large Language Model Based Autonomous Agents,'' \textit{Front. Comput. Sci.}, vol. 18, no. 6, p. 186345, 2024. doi: 10.1007/s11704-024-40231-1.

\bibitem{gao2023ragsurvey}
Y. Gao, Y. Xiong, X. Gao, K. Jia, J. Pan, Y. Bi, Y. Dai, J. Sun, and H. Wang, ``Retrieval-Augmented Generation for Large Language Models: A Survey,'' \textit{arXiv preprint arXiv:2312.10997}, 2023. doi: 10.48550/arXiv.2312.10997.

\bibitem{nist80030r1}
National Institute of Standards and Technology, ``Guide for Conducting Risk Assessments,'' \textit{NIST Special Publication (SP) 800-30 Rev. 1}, Gaithersburg, MD, 2012. doi: 10.6028/NIST.SP.800-30r1.

\bibitem{hong2024metagpt}
S. Hong, M. Zhuge, J. Chen, X. Zheng, C. Cheng, C. Zhang, J. Wang, Z. Wang, S. K. S. Yau, Z. Lin, L. Zhou, C. Ran, L. Xiao, C. Wu, and Q. Chen, ``MetaGPT: Meta Programming for A Multi-Agent Collaborative Framework,'' in \textit{Proc. Int. Conf. Learn. Represent. (ICLR 2024)}, 2024. [Online]. Available: https://openreview.net/forum?id=VtmBAGCN7o

\bibitem{wu2023autogen}
Q. Wu, G. Bansal, J. Zhang, Y. Wu, B. Li, E. Zhu, L. Jiang, X. Zhang, S. Zhang, J. Liu, A. H. Awadallah, R. W. White, D. Burger, and H. Wang, ``AutoGen: Enabling Next-Gen LLM Applications via Multi-Agent Conversation,'' \textit{arXiv preprint arXiv:2308.08155}, 2023. doi: 10.48550/arXiv.2308.08155.

\bibitem{boukari2021smishing}
B. E. Boukari, A. Ravi, and M. Msahli, ``Machine Learning Detection for SMiShing Frauds,'' in \textit{Proc. IEEE Consum. Commun. Netw. Conf. (CCNC)}, 2021, pp. 1--6. doi: 10.1109/CCNC49032.2021.9369640.

\bibitem{wooldridge2009multiagent}
M. Wooldridge, \textit{An Introduction to MultiAgent Systems}, 2nd ed. Chichester, U.K.: John Wiley \& Sons, 2009.

\bibitem{qian2024chatdev}
C. Qian, X. Cong, C. Yang, W. Chen, Y. Su, J. Xu, Z. Liu, and M. Sun, ``ChatDev: Communicative Agents for Software Development,'' in \textit{Proc. 62nd Annu. Meeting Assoc. Comput. Linguist. (ACL 2024)}, 2024, pp. 15174--15186. doi: 10.18653/v1/2024.acl-long.810.

\bibitem{lewis2020rag}
P. Lewis, E. Perez, A. Piktus, F. Petroni, V. Karpukhin, N. Goyal, H. K{\"u}ttler, M. Lewis, W.-t. Yih, T. Rockt{\"a}schel, S. Riedel, and D. Kiela, ``Retrieval-Augmented Generation for Knowledge-Intensive NLP Tasks,'' in \textit{Proc. Adv. Neural Inf. Process. Syst. (NeurIPS 2020)}, vol. 33, 2020, pp. 9459--9474.

\bibitem{reimers2019sbert}
N. Reimers and I. Gurevych, ``Sentence-BERT: Sentence Embeddings using Siamese BERT-Networks,'' in \textit{Proc. 2019 Conf. Empirical Methods Nat. Lang. Process. (EMNLP-IJCNLP)}, 2019, pp. 3982--3992. doi: 10.18653/v1/D19-1410.

\bibitem{malkov2020hnsw}
Y. A. Malkov and D. A. Yashunin, ``Efficient and Robust Approximate Nearest Neighbor Search Using Hierarchical Navigable Small World Graphs,'' \textit{IEEE Trans. Pattern Anal. Mach. Intell.}, vol. 42, no. 4, pp. 824--836, 2020. doi: 10.1109/TPAMI.2018.2889473.

\bibitem{munir2024transitrag}
M. B. Munir, Y. Cai, L. Khan, and B. Thuraisingham, ``Leveraging Multimodal Retrieval-Augmented Generation for Cyber Attack Detection,'' in \textit{Proc. IEEE 6th Int. Conf. Trust, Privacy \& Secur. Intell. Syst. Appl. (TPS-ISA)}, 2024, pp. 288--297. doi: 10.1109/tps-isa62245.2024.00046.

\bibitem{gao2023threatkg}
P. Gao, X. Liu, E. Choi, B. Hu, Z. Zheng, F. Shao, and M. Zitnik, ``ThreatKG: An AI-Powered System for Automated Open-Source Cyber Threat Intelligence Gathering,'' in \textit{Proc. 1st ACM Workshop Large AI Syst. Models (LEAM '23)}, 2023, pp. 1--8. doi: 10.1145/3627535.3638464.

\bibitem{rbi2021security}
Reserve Bank of India, ``Master Direction on Digital Payment Security Controls,'' \textit{Reserve Bank of India Notification RBI/2020-21/74}, Mumbai, India, Feb. 2021. [Online]. Available: https://www.rbi.org.in

\bibitem{radford2023whisper}
A. Radford, J. W. Kim, T. Xu, G. Brockman, C. McLeavey, and I. Sutskever, ``Robust Speech Recognition via Large-Scale Weak Supervision,'' in \textit{Proc. 40th Int. Conf. Mach. Learn. (ICML 2023)}, PMLR 202, 2023, pp. 28492--28518.

\bibitem{baek2019craft}
Y. Baek, B. Lee, D. Han, S. Yun, and H. Lee, ``Character Region Awareness for Text Detection,'' in \textit{Proc. IEEE/CVF Conf. Comput. Vis. Pattern Recognit. (CVPR)}, 2019, pp. 9365--9374. doi: 10.1109/CVPR.2019.00959.

\bibitem{shi2017crnn}
B. Shi, X. Bai, and C. Yao, ``An End-to-End Trainable Neural Network for Image-Based Sequence Recognition and Its Application to Scene Text Recognition,'' \textit{IEEE Trans. Pattern Anal. Mach. Intell.}, vol. 39, no. 11, pp. 2298--2304, 2017. doi: 10.1109/TPAMI.2016.2646371.

\bibitem{jung2022aasist}
J. Jung, H.-S. Heo, H. Tak, H.-j. Shim, J. S. Chung, B.-J. Lee, H.-J. Yu, and N. Evans, ``AASIST: Audio Anti-Spoofing Using Integrated Spectro-Temporal Graph Attention,'' in \textit{Proc. IEEE Int. Conf. Acoust., Speech Signal Process. (ICASSP)}, 2022, pp. 6367--6371. doi: 10.1109/ICASSP43922.2022.9747766.

\bibitem{wilson1927score}
E. B. Wilson, ``Probable Inference, the Law of Succession, and Statistical Inference,'' \textit{J. Am. Stat. Assoc.}, vol. 22, no. 158, pp. 209--212, 1927. doi: 10.1080/01621459.1927.10502953.

\end{thebibliography}

\end{document}
